\begin{abstract}
Prompt compression for privacy asks three questions: which content is sensitive, which is
relevant to the question, and which can be dropped to save tokens. This paper argues that a
fourth question decides whether the answers to the other three are usable. Once the identifier
has been removed or replaced, can the reader still tell which entity a surviving value belongs
to? We call that property \emph{referential integrity} and treat it as the object of design
rather than a side effect. On $2{,}728$ cases it factors into three conditions---the answer
survives, distinct entities do not collapse onto one surface token, and the question names the
same entity as the context. Two of the three are zeroed by exactly one configuration each, and
both times the task is lost ($0.3132$ and $0.3142$ F1 against $0.4937$); the third never reaches
zero and is reported as partial.

\textsc{DARK} preserves the property without keeping the value. Each unit is routed on
(sensitivity, relevance) to delete, alias, refill or keep, and the same alias map is written
into the question so that both halves keep a referent. The rule is training-free, reads
relevance from the question alone, and never consults the reference answer. Figure~1 states the
problem, the rule and the result on one page.

At a $50\%$ budget \textsc{DARK} reaches F1 $0.4937$ against $0.4831$ for the strongest
query-aware compressor, at $0\%$ recoverable against that method's $98.1\%$. The aggregate hides which half is which. On the $1{,}228$
questions that name a record identifier, \textsc{DARK} leads every compressor by $+0.1196$ to
$+0.4198$ and matches the uncompressed prompt at half its length; on the $1{,}500$ that name a
protected account, every compressor above it leaks on $89.5$--$100\%$ of those cases.
The zero is not an artifact of the pattern set: a model asked to recover the values returns
$36.4\%$ of them from the uncompressed prompt and none from \textsc{DARK}. Aliasing is what achieves that, and
aliasing alone does not make a compressor competitive---Quito-X on the aliased passage costs
$0.0371$ F1---which is why the routing decision, not the compressor, is the contribution.
Across budgets it is the only arm that improves as the budget grows without ever leaking; below
$r\approx0.10$ the one arm above it forwards a protected value in $21\%$ of cases.
Crossing the two design choices shows the operation is worth two to three times the region: it is the
contribution, and the region is its precondition. Code, prompts and
per-case records for every table and figure are at
\url{https://github.com/firesaku/PrivTrim-IoT}.
\end{abstract}
